\documentclass[11pt]{article}
\usepackage[margin=1in]{geometry}
\usepackage{natbib}
\usepackage{booktabs}
\usepackage{amsmath}
\usepackage{graphicx}
\usepackage{hyperref}
\usepackage{lineno}

\title{Disseminating cells in human oral tumours acquire an EMT cancer stem cell state that is predictive of metastasis}
\author{
Biddle, A.$^{1,*}$ \and Gammon, L.$^{1}$ \and Fazil, B.$^{1}$ \and Mackenzie, I.C.$^{1}$\\[6pt]
\small $^{1}$Centre for Oral Immunobiology and Regenerative Medicine, Institute of Dentistry,\\
\small Barts and The London School of Medicine and Dentistry, Queen Mary University of London, UK\\
\small $^{*}$Corresponding author: \texttt{a.biddle@qmul.ac.uk}
}
\date{}

\begin{document}
\maketitle

\begin{abstract}
Cancer stem cells (CSCs) undergo epithelial-mesenchymal transition (EMT) to drive metastatic dissemination in experimental cancer models. However, tumour cells undergoing EMT have not been observed disseminating into the tissue surrounding human tumour specimens, leaving the relevance to human cancer uncertain. We have previously identified both EpCAM and CD24 as markers of EMT CSCs with enhanced plasticity. This afforded the opportunity to investigate whether retention of EpCAM and CD24 alongside upregulation of the EMT marker Vimentin can identify disseminating EMT CSCs in human tumours. Examining disseminating tumour cells in over 12,000 imaging fields from 84 human oral cancer specimens, we see a significant enrichment of single EpCAM, CD24 and Vimentin co-stained cells disseminating beyond the tumour body in metastatic specimens. Through training an artificial neural network, these predict metastasis with high accuracy (cross-validated accuracy of 87--89\%). In this study, we have observed single disseminating EMT CSCs in human oral cancer specimens, and these are highly predictive of metastatic disease.
\end{abstract}

\section{Introduction}

In multiple types of carcinoma, cancer stem cells (CSCs) undergo epithelial-mesenchymal transition (EMT) to enable metastatic dissemination from the primary tumour \citep{biddle2011, liu2014, lawson2015}. This model of metastatic dissemination has been built from studies using murine models and human cancer cell line models. However, this process has not been observed in human tumours in the in vivo setting, leading to uncertainty over the relevance of these findings to human tumour metastasis \citep{bill2015, williams2019}.

A key complication with efforts to study metastatic processes in human tumours is the inability to trace cell lineage. As cancer cells exiting the tumour downregulate epithelial markers whilst undergoing EMT, they become indistinguishable from the mesenchymal non-tumour cells surrounding the tumour \citep{li2016}. This has represented a fundamental barrier to observing EMT-mediated dissemination in intact human tumour tissue.

Previous work has established that CSCs in oral squamous cell carcinoma (OSCC) exist in two distinct phenotypes: a proliferative epithelial state and a migratory EMT state \citep{biddle2011}. These EMT CSCs demonstrate enhanced plasticity, allowing them to switch between states and seed metastases through mesenchymal-epithelial transition (MET) at distant sites \citep{biddle2016, ocana2012, tsai2012}. Critically, we have identified that while epithelial keratins are lost during EMT, the epithelial cell adhesion molecule EpCAM is retained in EMT CSCs, along with CD24, a marker of the plastic EMT CSC sub-population \citep{biddle2016}. CD44 cannot be used as an EMT marker in the context of intact tissue as it requires trypsin degradation to yield differential expression \citep{biddle2013, mack2008}. Vimentin, on the other hand, accurately distinguishes EMT from epithelial tumour cells in immunofluorescent staining protocols \citep{biddle2016}.

EpCAM is a specific epithelial marker that is not expressed in stromal or immune cells---it is expressed exclusively in epithelia and epithelial-derived tumours \citep{keller2019}. This specificity, combined with Vimentin as a mesenchymal marker and CD24 as a plasticity marker, provides an opportunity to identify tumour cells that have undergone EMT and disseminated into the surrounding stromal region.

OSCC are an important health burden and one of the top ten cancers worldwide, with over 300,000 new cases per year. There is frequent metastatic spread to the lymph nodes of the neck; this is the single most important predictor of outcome and an important factor in treatment decisions \citep{sano2007}. In this study, we investigate whether EpCAM$^+$Vim$^+$CD24$^+$ cells can be detected disseminating beyond the tumour body in human OSCC specimens, and whether their presence is predictive of metastatic disease.

\section{Results}

\subsection{EpCAM, Vimentin and CD24 co-staining identifies the EMT stem cell state}

By combining EpCAM as a tumour lineage and EMT CSC marker, Vimentin as a mesenchymal marker, and CD24 as a plastic EMT CSC marker, we aimed to identify tumour cells that have undergone EMT and disseminated into the surrounding stromal region. We developed a protocol for automated 4-colour (3 markers + nuclear stain) immunofluorescent imaging and analysis of entire histopathological slide specimens, to test for co-localisation of the 3 markers in each individual cell across each specimen.

To determine whether this marker combination specifically identifies EMT CSCs, we stained the parental CA1 oral cancer cell line and an EMT-stem CA1 sub-line. EpCAM$^+$Vim$^+$CD24$^+$ cells were greatly enriched in the EMT-stem sub-line, comprising 41\% of the population, compared to 2.1\% in the CA1 line. Cells with this staining profile were absent from normal keratinocyte culture and cancer-associated fibroblast culture. There was very little Pan-keratin$^+$Vim$^+$CD24$^+$ staining, and no enrichment for Pan-keratin$^+$Vim$^+$CD24$^+$ cells in the EMT-stem sub-line.

Therefore, whilst epithelial keratins are lost, EpCAM is retained in cells undergoing EMT and an EpCAM$^+$Vim$^+$CD24$^+$ staining profile can be used as a marker for EMT CSCs in immunofluorescent staining protocols.

\subsection{EpCAM$^+$Vim$^+$CD24$^+$ cells are enriched in the stroma surrounding metastatic tumours}

Imaging the tumour body and adjacent stroma in sections of human OSCC specimens, we detected single cells co-expressing EpCAM, Vimentin and CD24 in the stromal region surrounding the tumour, confirming that these cells can be detected in human tumour specimens.

We next stratified 24 human primary OSCC specimens into 12 tumours that had evidence of lymph node metastasis or perineural spread, and 12 that remained metastasis-free, and stained them for EpCAM, Vimentin and CD24. Single cells co-expressing EpCAM, Vimentin and CD24 were abundant in the stroma surrounding metastatic tumours. This was not the case in non-metastatic tumours or normal epithelial regions. In contrast to EpCAM, pan-keratin staining did not identify cells in the stroma surrounding metastatic tumours.

Grey level intensities for EpCAM, Vimentin and CD24 were obtained for every nucleated cell in each imaging field. Expression of EpCAM, Vimentin and CD24 was then analysed for every nucleated cell in every imaging field that included both tumour and stroma (3,500 manually curated imaging fields across the 24 tumours containing 2,640,000 total nucleated cells).

EpCAM$^+$Vim$^+$CD24$^+$ cells were enriched in the stroma compared to the tumour body, and there was a much greater accumulation of EpCAM$^+$Vim$^+$CD24$^+$ cells in the stroma of metastatic tumours compared to non-metastatic tumours. Interestingly, this was not the case for EpCAM$^+$Vim$^+$CD24$^-$ cells, which were also enriched in the stroma but showed no difference between metastatic and non-metastatic tumours. Pan-keratin$^+$Vim$^+$CD24$^+$ cells were not detected.

\subsection{Validation in an independent cohort of 60 tumours}

To extend this analysis, we stained and imaged a further 60 tumours, evenly stratified on the same criteria. These displayed the same evidence of individual disseminating cells co-expressing EpCAM, Vimentin and CD24 in metastatic tumours only. The proportion of EpCAM$^+$Vim$^+$CD24$^+$ and EpCAM$^+$Vim$^+$CD24$^-$ cells was quantified for each cell in over 9,000 imaging fields at the tumour-stroma boundary. Consistent with the previous set of tumours, only EpCAM$^+$Vim$^+$CD24$^+$ cells were specifically enriched in the stroma of metastatic tumours.

\subsection{Confirmation using single-cell RNA sequencing data}

To explore whether these EpCAM$^+$Vim$^+$CD24$^+$ cells in the stroma may be non-tumour cell types, we analysed a published scRNAseq dataset for human head and neck cancer \citep{puram2017}. Analysing this dataset for EpCAM, Vimentin and CD24 co-expression, we found that 12\% of tumour cells (267/2,215) were EpCAM$^+$Vim$^+$CD24$^+$. In the non-tumour cells, only 0.8\% (29/3,687) were EpCAM$^+$Vim$^+$CD24$^+$. Therefore, the observed EpCAM$^+$Vim$^+$CD24$^+$ cells in our tumour specimens are highly likely to be a tumour cell population.

These findings demonstrate that an EpCAM$^+$Vim$^+$CD24$^+$ staining profile marks tumour cells disseminating into the surrounding stroma. The presence of disseminating tumour cells that express EpCAM but not CD24 did not correlate with metastasis. This highlights a requirement for the plasticity marker CD24 when identifying disseminating metastatic CSCs.

\subsection{Identification of EpCAM$^+$CD24$^+$Vim$^+$ CSCs enables clinical prediction using a machine learning approach}

We sought to determine whether the EpCAM$^+$CD24$^+$Vim$^+$ staining pattern could be predictive of metastasis. Starting with the EpCAM, Vimentin and CD24 immunofluorescence grey levels for each nucleated cell, we used a supervised machine learning approach to predict whether an imaging field comes from a metastatic or non-metastatic tumour. As a benchmark we used the pan-keratin, Vimentin and CD24 immunofluorescence grey levels, as we hypothesised that pan-keratin would provide an inferior predictive value than EpCAM given the absence of pan-keratin-expressing cells disseminating in the stroma.

We compared the performance accuracy (10-fold cross-validated F-score) of different machine learning classification algorithms \citep{pedregosa2011}. The best performing classifiers for EpCAM, Vimentin and CD24 were the artificial neural network (ANN) \citep{rumelhart1986} and support vector machine (SVM) \citep{smola2004}, with F1 accuracy scores of 91\% and 87\%, respectively. For the ANN, the area under the curve (AUC) accuracy score was 87\%, with 94\% sensitivity and 82\% specificity. Training with Pan-keratin, Vimentin and CD24 gave much worse prediction across all classifiers. Other classifiers tested included Naive Bayes \citep{zhang2005}, Random Forest \citep{dumont2009}, and k-Nearest Neighbours \citep{bentley1975}.

\subsection{Independent validation of predictive accuracy}

To extend this analysis of utility for metastasis prediction, we stained and imaged a further 60 tumour specimens (tumour batch 2). Over 9,000 imaging fields at the tumour-stroma boundary from 60 evenly stratified tumour specimens, containing over 8.5 million nucleated cells, were fed into an artificial neural network machine learning task. For this task, we recorded the predictive accuracy from the training and validation sets after each training epoch, which showed good alignment and an 89\% accuracy score after 12 training epochs.

To our knowledge, this is the first time immunofluorescent staining of human tumour tissue specimens has been combined with machine learning for the prediction of cancer metastasis. Previous studies using cytokeratin immunohistochemistry, clinicopathological data and serum biomarkers for clinical prediction via machine learning have achieved AUCs of 75\% in breast cancer \citep{tseng2019}, 80\% in OSCC \citep{bur2019}, and 82\% in colorectal cancer \citep{takamatsu2019}. Our approach achieves superior accuracy by specifically targeting the disseminating EMT CSC population.

\section{Discussion}

In this study, we have demonstrated that single disseminating tumour cells undergoing EMT can be identified in human oral cancer specimens through the co-expression of EpCAM, Vimentin and CD24. These cells are significantly enriched in the stroma surrounding metastatic tumours and can predict metastatic status with high accuracy using a machine learning approach.

The central finding is that the EMT CSC model, extensively validated in experimental systems \citep{biddle2011, liu2014, lawson2015}, is directly relevant to human cancer in vivo. By leveraging the retention of EpCAM during EMT---while epithelial keratins are lost---we overcame the long-standing challenge of distinguishing disseminating tumour cells from surrounding mesenchymal stromal cells \citep{li2016, bill2015}. The critical role of CD24 as a plasticity marker is highlighted by the finding that EpCAM$^+$Vim$^+$CD24$^-$ cells, while also enriched in the stroma, showed no association with metastasis. This is consistent with the model that not all EMT cells are equally capable of metastasis; rather, it is the plastic EMT CSC sub-population, marked by CD24, that drives metastatic dissemination \citep{biddle2016}.

The clinical utility of these findings is demonstrated by the high predictive accuracy for metastasis achieved through machine learning. The ANN classifier trained on EpCAM, Vimentin and CD24 grey level data achieved F1 scores of 87--91\% and an AUC of 87\%, outperforming previous machine learning approaches using clinicopathological data \citep{tseng2019, bur2019, takamatsu2019}. The independent validation in a second cohort of 60 tumours (89\% accuracy) provides strong evidence for the robustness and clinical potential of this approach.

The observation that pan-keratin staining fails to identify disseminating tumour cells underscores the importance of marker selection in EMT research \citep{tan2014, desousa2015}. Standard epithelial markers such as cytokeratins are lost during EMT, rendering them ineffective for tracking tumour cells in the mesenchymal compartment. EpCAM, with its retained expression during EMT \citep{keller2019}, provides a critical advantage for identifying these clinically important cells \citep{blancas2014}.

These findings have implications for clinical management of OSCC. Current methods for predicting lymph node metastasis have limited accuracy, yet metastatic spread is the single most important prognostic factor \citep{sano2007}. An immunofluorescence-based assay combined with machine learning analysis could provide a more accurate predictive tool, potentially guiding treatment decisions regarding elective neck dissection.

\section{Materials and Methods}

\subsection{Cell lines and culture}

The CA1 oral squamous cell carcinoma cell line and its EMT-stem sub-line were used for in vitro validation. Cancer-associated fibroblasts and N/TERT keratinocytes \citep{smits2017} were used as negative controls.

\subsection{Human tumour specimens}

Eighty-four human primary OSCC specimens were analysed in two batches. Batch 1 comprised 24 tumours (12 metastatic, 12 non-metastatic), and batch 2 comprised 60 tumours (30 metastatic, 30 non-metastatic), stratified by evidence of lymph node metastasis or perineural spread.

\subsection{Immunofluorescent staining and imaging}

Four-colour immunofluorescent staining was performed for EpCAM (yellow), Vimentin (red), and CD24 (green) with DAPI nuclear stain (blue). Parallel sections were stained with pan-keratin replacing EpCAM. Automated imaging was performed across entire histopathological slide specimens.

\subsection{Image analysis and segmentation}

An image segmentation protocol was developed to separate the tumour body from the adjacent stroma. Grey level intensities for EpCAM, Vimentin and CD24 were obtained for every nucleated cell in each imaging field. In batch 1, 3,500 manually curated imaging fields containing 2,640,000 nucleated cells were analysed. In batch 2, over 9,000 imaging fields containing over 8.5 million nucleated cells were analysed.

\subsection{Machine learning}

Supervised machine learning was performed using scikit-learn \citep{pedregosa2011}. Classification algorithms compared included: artificial neural network \citep{rumelhart1986}, support vector machine \citep{smola2004}, Naive Bayes \citep{zhang2005}, Random Forest \citep{dumont2009}, and k-Nearest Neighbours \citep{bentley1975}. Performance was assessed using 10-fold cross-validated F1 scores. For batch 2, an ANN classifier was trained and tested independently, with accuracy and loss tracked over 14 training epochs.

\subsection{scRNAseq analysis}

A published single-cell RNA sequencing dataset for human head and neck squamous cell carcinoma \citep{puram2017} was analysed for co-expression of EpCAM, Vimentin and CD24 in tumour and non-tumour cell populations.

\bibliographystyle{plainnat}
\bibliography{references}

\end{document}
